\documentclass[10pt,a4paper]{article}

%double si on veut une version prof et une version élève. simple si on veut une seule version.
\newcommand{\buildmode}{double}

\InputIfFileExists{version.tex}{}{}
% Version (définie par le script)
\providecommand{\version}{eleve}

\usepackage{../feuilleexo}

%%%à modifier à chaque fois

\renewcommand{\niveau}{1G}
\renewcommand{\chapitre}{Fonctions polynômes du second degré}
\renewcommand{\numchap}{1}
\renewcommand{\theme}{Fonctions}

\begin{document}

%\legendeicones

\vspace{-0.3cm}

\begin{prerequis}
\begin{itemize}
    \item Développer et factoriser une expression algébrique (double distributivité, identités remarquables).
    \item Résoudre une équation du premier degré et un produit nul.
    \item Connaître le vocabulaire d'une fonction : image, antécédent, sens de variation, extremum.
    \item Lire et compléter un tableau de valeurs, placer des points dans un repère.
\end{itemize}
\end{prerequis}

\begin{multicols}{2}

\section{Formes d'une fonction polynôme du second degré}

\begin{exo}
On considère $f(x) = 2x^2 - 12x + 7$.
\begin{enumerate}[leftmargin=*]
    \item Donner les coefficients $a$, $b$ et $c$ de la forme développée.
    \item Vérifier que la forme canonique de $f$ est $f(x) = 2(x-3)^2 - 11$.
    \item En déduire la forme factorisée de $f$, si elle existe.
\end{enumerate}
\end{exo}

\begin{exo}
Développer puis réduire les expressions suivantes.
\begin{tasks}(2)
\task $3(x-1)^2+4$
\task $-2(x+5)^2-1$
\task $(x-2)(3x+7)$
\task $\dfrac12(x+4)^2-6$
\end{tasks}
\end{exo}

\begin{exo}[Compléter les étapes][appro]
Soit $g(x) = x^2 + 6x - 1$. On veut écrire $g$ sous forme canonique.
\begin{leftbar}
$g(x) = x^2 + 6x - 1$

$\phantom{g(x)} = (x + \dots)^2 - \dots - 1$

$\phantom{g(x)} = (x + \dots)^2 - \dots$
\end{leftbar}
Compléter les pointillés, puis donner le sommet de la parabole représentant $g$.
\end{exo}

\section{Équations et signe du trinôme}

\begin{exo}
Résoudre dans $\R$ les équations suivantes.
\begin{tasks}(2)
\task $x^2-5x+6=0$
\task $2x^2+3x+5=0$
\task $-x^2+4x-4=0$
\task $x^2-7=0$
\end{tasks}
\end{exo}

\begin{exo}
Étudier le signe de $h(x) = -3x^2+2x+1$ sur $\R$, et en déduire les solutions de l'inéquation $h(x) \leqslant 0$.
\end{exo}

\begin{exo}[][appro]
Un élève affirme :
\begin{leftbar}
\og Une fonction du second degré est toujours positive, sauf entre ses deux racines quand elle en a. \fg{}
\end{leftbar}
Cette affirmation est-elle correcte ? Justifier à l'aide d'un contre-exemple si nécessaire.
\end{exo}

\begin{exo}[][bac]
Une entreprise fabrique $x$ objets par jour ($x \in [0;40]$). Son bénéfice journalier, en euros, est modélisé par $B(x) = -2x^2+80x-150$.
\begin{enumerate}[leftmargin=*]
    \item Déterminer le nombre d'objets à produire pour que l'entreprise ne soit pas déficitaire.
    \item Déterminer la production journalière qui maximise le bénéfice, et calculer ce bénéfice maximal.
\end{enumerate}
\end{exo}

\section{Variations, extremum et représentation graphique}

\begin{exo}
Pour chacune des fonctions suivantes, donner la forme canonique, le sens de variation et dresser le tableau de variations complet.
\begin{enumerate}
    \item $f(x) = x^2 - 4x + 1$
    \item $g(x) = -2x^2 + 8x - 3$
\end{enumerate}
\end{exo}

\begin{exo}[][appro]
Déterminer l'expression de la fonction polynôme du second degré $f$ dont la parabole représentative admet pour sommet $S(2\,;\,-3)$ et passe par le point $A(0\,;\,5)$.
\end{exo}

\begin{exo}[Lecture graphique][appro]
On donne ci-dessous la représentation graphique $\mathcal{C}_f$ d'une fonction polynôme du second degré $f$.

\begin{center}
\begin{tikzpicture}
\begin{axis}[
    width=6.5cm, height=5.5cm,
    xmin=-2.5, xmax=4.5, ymin=-5.5, ymax=2.5,
    axis lines=middle,
    xlabel=$x$, ylabel=$y$,
    xtick={-2,-1,1,2,3,4}, ytick={-5,-4,-3,-2,-1,1,2},
    ticklabel style={font=\tiny},
    label style={font=\small},
    axis line style={-Latex},
]
\addplot[domain=-2.2:4.2, samples=100, thick] {(x-1)*(x-3)-1};
\end{axis}
\end{tikzpicture}
\end{center}

\begin{enumerate}[leftmargin=*]
    \item Par lecture graphique, donner le signe de $f$ sur $[-2{,}5\,;\,4{,}5]$ ainsi que les coordonnées du sommet de $\mathcal{C}_f$.
    \item En déduire la forme canonique de $f$, puis vérifier sa forme développée en sachant que $f(0)=2$.
\end{enumerate}
\end{exo}

\section{Problèmes de synthèse}

\begin{exo}[Aire maximale][bac]
On dispose d'un fil de fer de longueur $20$ cm que l'on plie pour former un rectangle de largeur $x$ (en cm), avec $x \in [0\,;\,10]$.
\begin{enumerate}[leftmargin=*]
    \item Exprimer, en fonction de $x$, la longueur du rectangle puis son aire $A(x)$.
    \item Montrer que $A$ admet un maximum que l'on précisera, ainsi que la valeur de $x$ pour laquelle il est atteint.
    \item Quelle est la nature du rectangle d'aire maximale ?
\end{enumerate}
\end{exo}

\begin{exo}[Coût de production][bac]
Le coût de production de $x$ centaines d'articles, en milliers d'euros, est donné par $C(x) = 0{,}5x^2 - 3x + 8$ pour $x \in [0\,;\,8]$. Chaque centaine d'articles est vendue $2$ milliers d'euros.
\begin{enumerate}[leftmargin=*]
    \item Exprimer le bénéfice $B(x)$ en fonction de $x$.
    \item Déterminer pour quelles valeurs de $x$ l'entreprise réalise un bénéfice.
    \item Déterminer la quantité $x$ qui maximise le bénéfice, et la valeur de ce bénéfice.
\end{enumerate}
\end{exo}

\end{multicols}

\end{document}
